\section{Conclusion}

\label{sec:conclusion}
% ============================================================

Sovereign DeFi establishes that cross-chain bridge governance need not be centralized, compliance need not be advisory, and sovereignty need not compromise security. Through the separation of security, economic, and protocol authority domains, Lux Teleport enables each connected chain to govern its own bridge instance while sharing the security of a unified MPC threshold signature network.

The three contributions---per-chain governance with proven isolation, protocol-level Shariah compliance with SAB authority, and ERC-3643-compatible securities bridging---address distinct regulatory and community requirements through a unified architecture. The formal proofs (governance isolation, exit guarantee, fee cap invariant) provide the assurances that institutional and regulatory participants require.

The fee competition game between sovereign instances produces efficient equilibria well below the hard cap, while the xLUX flywheel demonstrates how fee aggregation creates sustainable economic incentives. The white-label architecture transforms Teleport from a single bridge into a bridge \emph{protocol} that any chain can instantiate with its own governance and branding.

The immediate impact is concrete: protocol-level Shariah enforcement opens DeFi to 1.8~billion Muslims who are currently excluded by interest-bearing yield mechanisms, securities compliance preservation enables regulated assets to move cross-chain without violating jurisdictional constraints, and sovereign governance ensures that no external party can extract rent from a chain's bridge traffic.

All contracts, classification data, and the settlement registry are open-source and auditable on-chain.

% ============================================================
% References
% ============================================================
\begin{thebibliography}{99}

\bibitem{frost} C.~Komlo and I.~Goldberg, ``FROST: Flexible Round-Optimized Schnorr Threshold Signatures,'' in \emph{Selected Areas in Cryptography (SAC)}, 2020.

\bibitem{cggmp21} R.~Canetti, R.~Gennaro, S.~Goldfeder, N.~Makriyannis, and U.~Peled, ``UC Non-Interactive, Proactive, Threshold ECDSA with Identifiable Aborts,'' in \emph{ACM CCS}, 2020.

\bibitem{wormhole} Wormhole Foundation, ``Wormhole: A Generic Message Passing Protocol,'' Whitepaper, 2022.

\bibitem{layerzero} R.~Zarick, B.~Pellegrino, and C.~Banister, ``LayerZero: An Omnichain Interoperability Protocol,'' Whitepaper, 2022.

\bibitem{axelar} S.~Layr \emph{et al.}, ``Axelar: Connecting Web3,'' Whitepaper, 2022.

\bibitem{ibc} C.~Goes \emph{et al.}, ``The Interblockchain Communication Protocol: An Overview,'' Cosmos Documentation, 2020.

\bibitem{tbtc} M.~Luongo and C.~Ream, ``tBTC: A Decentralized Redeemable BTC-backed ERC-20 Token,'' Keep Network, 2020.

\bibitem{erc3643} J.~Music, D.~Music, L.~Music, and R.~Music, ``ERC-3643: T-REX Token for Regulated EXchanges,'' Ethereum Improvement Proposal, 2021.

\bibitem{usmani2002} M.~T.~Usmani, \emph{An Introduction to Islamic Finance}, Kluwer Law International, 2002.

\bibitem{iqbal2011} Z.~Iqbal and A.~Mirakhor, \emph{An Introduction to Islamic Finance: Theory and Practice}, 2nd ed., Wiley, 2011.

\bibitem{elgari2003} M.~A.~El-Gamal, ```Interest' and the Paradox of Contemporary Islamic Law and Finance,'' \emph{Fordham International Law Journal}, vol.~27, no.~1, 2003.

\bibitem{aaoifi2015} AAOIFI, \emph{Shari'ah Standards}, Accounting and Auditing Organization for Islamic Financial Institutions, Bahrain, 2015.

\bibitem{ifdi2024} ICD-REFINITIV, ``Islamic Finance Development Indicator Report,'' 2024.

\bibitem{worldbank2022} World Bank, ``Global Findex Database 2021: Financial Inclusion, Digital Payments, and Resilience in the Age of COVID-19,'' 2022.

\bibitem{mica2023} European Parliament, ``Regulation (EU) 2023/1114 on Markets in Crypto-Assets (MiCA),'' \emph{Official Journal of the European Union}, 2023.

\bibitem{polymath2018} Polymath Network, ``The Securities Token Standard (ST-20),'' Whitepaper, 2018.

\bibitem{mrhool2022} A.~Mrhool and S.~Al-Bassam, ``Shariah Compliance in Decentralized Finance: A Systematic Review,'' \emph{International Journal of Islamic and Middle Eastern Finance and Management}, vol.~15, no.~3, 2022.

\bibitem{shariachain2021} F.~Ullah, I.~Mehmood, and A.~Khan, ``ShariaChain: A Blockchain Framework for Shariah-Compliant Financial Services,'' in \emph{IEEE International Conference on Blockchain}, 2021.

\bibitem{haqq2023} HAQQ Network, ``Islamic Coin: Shariah-Compliant Digital Money,'' Whitepaper, 2023.

\bibitem{roughgarden2021} T.~Roughgarden, ``Transaction Fee Mechanism Design,'' in \emph{ACM Conference on Economics and Computation (EC)}, 2021.

\bibitem{rekt} Rekt News, ``Leaderboard,'' \url{https://rekt.news/leaderboard/}, 2024.

\bibitem{erc4626} J.~Bebis \emph{et al.}, ``EIP-4626: Tokenized Vault Standard,'' Ethereum Improvement Proposal, 2022.

\end{thebibliography}

